\section{Tesis: lo que te toca}
Documentos que valida el departamento en esta opción:
\begin{itemize}
  \item Acta de nacimiento
  \item CURP
  \item Carta de autorización o de excepción para publicar en el Repositorio Institucional
  \item Carta de presentación del pasante y asesor
  \item Certificado de estudios de bachillerato
  \item Certificado de estudios de licenciatura
  \item Certificado de servicio social
  \item Comprobante del sistema de Seguimiento de Egresados
  \item Constancia de autenticación de bachillerato
  \item Constancia de no adeudo a la Biblioteca Central
  \item Constancia de no adeudo a la Biblioteca de la Facultad
  \item Constancia de no adeudo a la Dirección de Control Escolar de la UAEM
  \item Constancia o certificado de aprobación del idioma inglés
  \item Contra-recibo de la donación de libros a la Biblioteca de la Facultad
  \item Dictamen de autorización de los HH. Consejos
  \item Formato de datos personales (para el título)
  \item Recibo de no adeudo a la UAEM
  \item Recibo de pago: derechos de evaluación profesional (\$3,024)
  \item Recibo de pago: expedición de título (\$1,187)
  \item Reporte de porcentaje de similitud (iThenticate)
  \item Solicitud de prórroga del tema (pasante y asesor)
  \item Solicitud de revisión firmada (pasante y asesor)
  \item Solicitud de sustentación firmada (pasante y asesor)
  \item Solicitud para presentar la evaluación profesional (trabajo escrito)
  \item Solicitud y protocolo firmados (pasante y asesor)
  \item Trabajo escrito en PDF para revisión
  \item Trabajo escrito final en PDF
  \item Voto aprobatorio (lo firman tu asesor y tus revisores)
\end{itemize}
Lo que registra el departamento:
\begin{itemize}
  \item Aprobación del protocolo y oficio de aprobación del tema
  \item Fecha del examen previo y sínodo
  \item Autorización de impresión
  \item Fecha del examen definitivo y jurado
  \item Registrar resultado
\end{itemize}
Renglones de la verificación final (FORMATO DE INTEGRACIÓN DE EXPEDIENTE: MODALIDADES CON TRABAJO ESCRITO):
\begin{enumerate}
  \item Dos fotografías tamaño semitítulo, blanco y negro, papel mate, fondo gris claro, de frente, con retoque, varones con camisa blanca, saco y corbata oscuros, mujeres con blazer oscuro
  \item Acta de evaluación profesional
  \item Copia del certificado de bachillerato con firma y sello de cotejo con el original
  \item Constancia de autenticación si el bachillerato no es UAEM \emph{(solo si aplica)}
  \item Dictamen de revalidación expedido por la Secretaría de Educación Pública (sólo si aplica) \emph{(solo si aplica)}
  \item Acta de nacimiento
  \item Certificado total de licenciatura
  \item Solicitud por escrito del egresado dirigida a la Subdirección Académica con firma y sello de la misma
  \item Constancia de no adeudo a Servicios Escolares
  \item Autorización para continuar con trámites de titulación con fundamento en el artículo 5° del Reglamento de Evaluación Profesional (sólo si aplica) \emph{(solo si aplica)}
  \item Formato de revisión para autorización de evaluación profesional
  \item Formato de integración de expediente
  \item Formato de datos del interesado para expedición de título profesional
\end{enumerate}
\subsection{La bandeja «Validar trámites»}
La bandeja reúne los trámites que esperan algo del departamento. Abre el folio para ver su expediente.
\captura{tesis_responsable_titulacion_01_bandeja}{Bandeja «Validar trámites»: los trámites con documentos por validar.}
\subsection{El expediente del trámite}
Cada documento aparece con su estado, el enlace a su PDF y los botones \boton{Validar} y \boton{Devolver}.
\captura{tesis_responsable_titulacion_02_expediente}{Expediente del trámite: cada documento con su PDF y los botones «Validar» y «Devolver».}
\subsection{Validar documentos}
Abre el PDF, revisa que sea el documento correcto y legible, y pulsa \boton{Validar}.
\captura{tesis_responsable_titulacion_03_validar}{Validar un documento: revisas el PDF y pulsas «Validar».}
\subsection{Devolver un documento}
Escribe en \emph{Observaciones} qué debe corregir el pasante y pulsa \boton{Devolver}. El pasante lo ve en su trámite y vuelve a subirlo.
\captura{tesis_responsable_titulacion_05_devolver}{Devolver un documento: escribes la observación y pulsas «Devolver»; el pasante la ve al entrar a su trámite.}
\subsection{Registrar lo que le toca al departamento}
Fechas, lugar, oficios y, al final, el resultado y el número de acta. Al registrarlo se completa la etapa.
\captura{tesis_responsable_titulacion_06_hito_sustentacion}{Registrar lo que le toca al departamento en esta etapa y completarla.}
\captura{tesis_responsable_titulacion_07_hito_resultado}{Registrar el resultado de la evaluación y el número de acta.}
\subsection{Visto bueno de la Jefatura}
Los oficios que llevan el visto bueno de la Jefatura esperan en \emph{Firmas pendientes}.
\captura{tesis_responsable_titulacion_04_firmas_jefatura}{La Jefatura de Titulación da su visto bueno en «Firmas pendientes».}
\subsection{Expedientes a Servicios Escolares}
Al concluir, el trámite pasa a esta bandeja para la verificación final.
\captura{tesis_responsable_titulacion_08_expedientes}{«Expedientes a Servicios Escolares»: los trámites concluidos por verificar y exportar.}
\subsection{Verificación final}
Revisa cada renglón del formato de integración y marca \emph{Entregado y revisado} o \emph{No aplica}. Después pulsa \boton{Guardar lista} y \boton{Hago constar que el expediente está completo y verificado}.
\captura{tesis_responsable_titulacion_09_verificacion}{Verificación final: cada renglón del formato de integración se marca como entregado o «No aplica».}
\subsection{Exportar para Servicios Escolares}
Con el expediente verificado, exporta la carpeta comprimida.
\captura{tesis_responsable_titulacion_10_exportar}{Expediente verificado: se exporta la carpeta comprimida para Servicios Escolares.}
